\documentclass[11pt]{article}
\usepackage[margin=1in]{geometry}
\usepackage{amsmath,amssymb,amsthm}
\usepackage{graphicx}
\usepackage[hidelinks]{hyperref}

\newtheorem{theorem}{Theorem}
\newtheorem{corollary}[theorem]{Corollary}
\newtheorem{proposition}[theorem]{Proposition}
\newcommand{\Lip}{\operatorname{Lip}}
\newcommand{\Lipz}{\operatorname{Lip}_0}
\newcommand{\R}{\mathbb R}

\title{Pointed Lipschitz fractal interpolation on intervals and curves}
\author{FA/Banach run \texttt{fa\_banach\_001}, lane 1}
\date{July 21, 2026}

\begin{document}
\maketitle

\section*{Source metadata and status}

The source is P.~R.~Massopust, \emph{Fractal Interpolation over Curves},
arXiv:2209.01033 (2022), published in \emph{Contemporary Mathematics} 797
(2024), 61--73 \cite{Massopust}.  Section~7, page~13 of the arXiv PDF, leaves
two research directions to the reader.  The second is to replace the unital
Lipschitz algebra by the algebra of bounded real-valued Lipschitz functions
vanishing at a base point.

\medskip
\noindent\textbf{Claimed status: full solution, likely valid.}
The theorem below carries out that pointed-Lipschitz direction for the interval
partitions used in fractal interpolation, gives exact junction conditions, and
transports the result to simple bi-Lipschitz curves in Banach spaces.

\section*{The source direction}

\begin{center}
\includegraphics[width=\textwidth]{figures/open_problem_crop.png}
\end{center}

In the notation of the source, the direction is to replace
$\Lip(X,F)$ by the pointed space
\[
  \Lipz(X)=\{f:X\to\R:f\text{ is Lipschitz and }f(x_0)=0\}.
\]
On a bounded pointed metric space this is a Banach space for the norm
$\|f\|_{\Lipz}=\Lip(f)$; the base-point condition implies
$\|f\|_\infty\leq \operatorname{diam}(X)\Lip(f)$.

\section*{Proof intuition}

The pointed condition changes the join-up problem.  At an internal knot, the
right branch sees $f(0)=0$, whereas the left branch sees the unrestricted value
$f(1)$.  Consequently, an RB formula acts on \emph{all} of $\Lipz$ precisely
when the left vertical multiplier vanishes at that endpoint and the two affine
terms agree there.

For contraction, multiplication contributes two terms:
$\|s\|_\infty\Lip(f-g)$ and $\Lip(s)\|f-g\|_\infty$.
The pointed estimate $\|f-g\|_\infty\leq\Lip(f-g)$ combines them, while
composition with an inverse partition branch contributes
$\Lip(h_i^{-1})$.  The junction conditions make the branchwise estimates glue
without any additional loss on an interval.

\section*{Pointed fixed-point theorem}

Fix a partition
\[
  0=t_0<t_1<\cdots<t_n=1.
\]
For $1\leq i\leq n$, let
$h_i:[0,1]\to[t_{i-1},t_i]$ be an increasing bi-Lipschitz homeomorphism,
with $h_i(0)=t_{i-1}$ and $h_i(1)=t_i$.  Put
\[
  \lambda_i=\Lip(h_i^{-1}),\qquad
  a_i=\|s_i\|_\infty,\qquad b_i=\Lip(s_i),
\]
where $q_i,s_i\in\Lip([0,1])$.

For $f\in\Lipz([0,1])$, consider the branch formula
\begin{equation}\label{eq:RB}
 (Tf)(x)=q_i(h_i^{-1}x)+s_i(h_i^{-1}x)f(h_i^{-1}x),
 \qquad x\in[t_{i-1},t_i].
\end{equation}

\begin{theorem}\label{thm:main}
The formula \eqref{eq:RB} defines an affine self-map of all of
$\Lipz([0,1])$ if and only if
\begin{equation}\label{eq:join}
 q_1(0)=0,\qquad
 s_i(1)=0,\quad q_i(1)=q_{i+1}(0)\qquad(1\leq i<n).
\end{equation}
Under these conditions,
\begin{equation}\label{eq:constant}
 \Lip(Tf-Tg)\leq c\,\Lip(f-g),
 \qquad
 c:=\max_{1\leq i\leq n}\lambda_i(a_i+b_i).
\end{equation}
In particular, if $c<1$, there is a unique
$\psi\in\Lipz([0,1])$ satisfying
\begin{equation}\label{eq:selfref}
 \psi(h_i(t))=q_i(t)+s_i(t)\psi(t),
 \qquad t\in[0,1],\quad 1\leq i\leq n.
\end{equation}
For every $f_0\in\Lipz([0,1])$, the iterates $f_k=T^k f_0$ converge to
$\psi$ in pointed Lipschitz norm and obey
\begin{equation}\label{eq:error}
 \Lip(\psi-f_k)\leq \frac{c^k}{1-c}\Lip(Tf_0-f_0).
\end{equation}
\end{theorem}

\begin{proof}
First suppose \eqref{eq:RB} defines a self-map on all of $\Lipz$.
At $0=h_1(0)$, taking $f=0$ gives $(Tf)(0)=q_1(0)$, hence
$q_1(0)=0$.  At the knot $t_i=h_i(1)=h_{i+1}(0)$, equality of the two
branch values gives, for every $f\in\Lipz$,
\[
 q_i(1)+s_i(1)f(1)=q_{i+1}(0)+s_{i+1}(0)f(0)=q_{i+1}(0).
\]
Taking first $f=0$ and then $f(t)=t$ proves the remaining conditions in
\eqref{eq:join}.

Conversely, assume \eqref{eq:join}.  It makes the two definitions in
\eqref{eq:RB} agree at every internal knot, and it gives $(Tf)(0)=0$.
On the $i$th interval, composition and the product estimate show that
\begin{align*}
 \Lip\big((Tf)|_{[t_{i-1},t_i]}\big)
 \leq \lambda_i\big(
   \Lip(q_i)+a_i\Lip(f)+b_i\|f\|_\infty
 \big),
\end{align*}
so $Tf$ is continuous and finitely piecewise Lipschitz.  Such a function on
an interval is globally Lipschitz: telescope across the intervening knots and
use that the subinterval lengths sum to the total length.  Thus $T$ maps
$\Lipz$ to itself.

Let $u=f-g$.  On the $i$th branch,
\[
 (Tf-Tg)(x)=(s_i u)(h_i^{-1}x),
\]
and hence
\begin{align*}
 \Lip\big((Tf-Tg)|_{[t_{i-1},t_i]}\big)
 &\leq\lambda_i\big(a_i\Lip(u)+b_i\|u\|_\infty\big)\\
 &\leq\lambda_i(a_i+b_i)\Lip(u),
\end{align*}
because $u(0)=0$ and the interval has diameter one.  Moreover, at every
internal knot the left difference is $s_i(1)u(1)=0$ and the right difference
is $s_{i+1}(0)u(0)=0$.  The same telescoping argument therefore gives the
global estimate \eqref{eq:constant}.

Finally, $\Lipz([0,1])$ is complete for the norm $\Lip(\cdot)$.  Indeed, a
Cauchy sequence in that norm is uniformly Cauchy because every member
vanishes at zero; its uniform limit is pointed Lipschitz, and passage to the
limit in difference quotients gives convergence in Lipschitz norm.  If $c<1$,
Banach's fixed-point theorem now gives the unique fixed point, convergence of
the iterates, and \eqref{eq:error}.  The fixed-point equation is exactly
\eqref{eq:selfref}.
\end{proof}

\section*{Transfer to curves}

\begin{corollary}\label{cor:curve}
Let $E$ be a Banach space and let
$\gamma:[0,1]\to E$ be a bi-Lipschitz homeomorphism onto its trace
$X=\gamma([0,1])$.  Set $x_0=\gamma(0)$ and
\[
 H_i=\gamma\circ h_i\circ\gamma^{-1},\qquad
 Q_i=q_i\circ\gamma^{-1},\qquad S_i=s_i\circ\gamma^{-1}.
\]
If \eqref{eq:join} and $c<1$ hold, there is a unique
$\Psi\in\Lipz(X)$ such that
\[
 \Psi(H_i x)=Q_i(x)+S_i(x)\Psi(x),
 \qquad x\in X,\quad 1\leq i\leq n.
\]
It is $\Psi=\psi\circ\gamma^{-1}$, where $\psi$ is the fixed point from
Theorem~\ref{thm:main}.  The transported RB iterates converge in
$\Lipz(X)$.
\end{corollary}

\begin{proof}
Composition with $\gamma$ and $\gamma^{-1}$ gives mutually inverse bounded
linear maps between $\Lipz(X)$ and $\Lipz([0,1])$.  Under these maps, the
curve RB operator is conjugate to \eqref{eq:RB}.  Existence, uniqueness, and
convergence therefore transfer directly.
\end{proof}

In particular, the corollary applies to every simple regular $C^1$ arc on a
compact interval once it is given a bi-Lipschitz parametrization (for example,
the usual regular embedded arcs used by interval-based fractal interpolation).

\section*{Why the contraction constant needs both terms}

The two summands in \eqref{eq:constant} cannot generally be replaced by their
maximum.  Take $n=1$, $h_1$ the identity, $q_1=0$, and
$s_1(t)=a t$ with $0<a<1$.  This defines $Tf(t)=atf(t)$ on all of
$\Lipz([0,1])$.  For $f(t)=t$,
\[
 \Lip(f)=1,
 \qquad
 \Lip(Tf)=\Lip(at^2)=2a.
\]
Thus $T$ is not a contraction whenever $a>1/2$, even though both
$\|s_1\|_\infty=a$ and $\Lip(s_1)=a$ are separately less than one.  The
coefficient in Theorem~\ref{thm:main} is $2a$, and is attained on this test
function.

\section*{Verification, scope, and novelty}

The proof is exact and has no numerical or computer-assisted component.  Its
checks are: completeness of the pointed norm; necessity and sufficiency of the
junction identities; the product Lipschitz estimate; inverse-branch expansion;
finite interval gluing; and conjugacy under a bi-Lipschitz curve
parametrization.

The result completes the source's pointed-Lipschitz direction for the natural
interval partition model and its simple bi-Lipschitz curve transports.  It does
not address the source's separate, much broader direction concerning arbitrary
Banach or Fr\'echet submanifolds.

The local run indexes were searched for arXiv:2209.01033, the title, pointed
Lipschitz algebras, $\Lipz$, and RB operators.  Bounded web searches on
July~21, 2026 used the exact Section~7 wording and the combinations
``pointed Lipschitz'', ``$\Lipz$'', ``fractal interpolation'', and
``Read--Bajraktarevi\'c''.  The arXiv record, 2024 publication metadata, and
current author-publication search results were checked.  No later work
explicitly carrying out this direction was found.  Novelty confidence is
moderate because the fixed-point argument is elementary; proof confidence is
high.

\medskip
\noindent\textbf{Human review recommendation.}
Review as a likely valid full answer to the second Section~7 direction.  The
most useful check is that ``replace by $\Lipz$'' is intended as a self-map on
the whole pointed algebra; under that literal interpretation, the junction
conditions in \eqref{eq:join} are forced, not optional.

\begin{thebibliography}{9}
\bibitem{Massopust}
P.~R.~Massopust,
\emph{Fractal Interpolation over Curves},
arXiv:2209.01033 (2022); \emph{Contemporary Mathematics} 797 (2024), 61--73,
doi:10.1090/conm/797/15934.
\end{thebibliography}

\end{document}
